\section{Experiments and validation}
\label{sec:experiments}

The experimental program follows the claim hierarchy. Exact correctness tests precede mechanism studies, and mechanism studies precede downstream benchmarks. This ordering prevents benchmark performance from masking a failure of the stated spectral construction.

\subsection{Gate A: correctness criteria}

The Gate-A contract contains 36 prespecified correctness criteria.  Its
executable suite covers strict graph validation; root, phase, and mapped-pole
geometry; agreement with an independently assembled dense oracle; exact
one-level and multilevel identities; adjoint and additive reconstruction;
weighted spectral identities; permutation and repeated-eigenspace cases;
finite-order error, frame, perturbation, locality, and operation-count bounds;
and counterexamples to target-specific long-range sensitivity.  Each row emits
typed numerical evidence and is cross-checked against its declared graph,
depth, root, and polynomial-degree fixtures.  Exact tests use double or
complex-double precision with row-specific tolerances or analytical bounds,
rather than a global threshold.  The current suite passes all 36 rows, but
Gate A remains closed pending a fresh independent semantic review of the final
public interfaces.  We therefore treat the suite as implementation validation,
not as independent evidence for the paper's empirical claims.

\subsection{Mechanism studies}

The first mechanism study compares a magnitude-only band objective with a complex response objective. On a seeded 600-node sphere graph, we fit eight roots from five initializations for each objective. The input combines equal-energy low-frequency and localized packet components. The retained aggregate reports recovery NMSE of $0.673\pm0.271$ versus $0.360\pm0.207$, leakage of $0.323\pm0.176$ versus $0.156\pm0.151$, and target-band energy retention of $0.666\pm0.099$ versus $0.834\pm0.011$ for magnitude and complex fitting, respectively. The dispersion is the reported standard deviation across initializations. \Cref{fig:sphere-decomposition} shows the graph, components, learned channels, spectral response, and root--pole trajectories. However, the artifact retains only these aggregates and one selected complex run, not the ten per-initialization records. The values therefore illustrate the intended distinction between energy selection and phase-sensitive recovery but are not yet claim-admissible or evidence of superiority to external methods.

\begin{figure*}[t]
\centering
\includegraphics[width=\textwidth]{fig_sphere_decomposition.pdf}
\caption{Diagnostic phase-sensitive recovery on a 600-node sphere kNN graph. (A) The symmetric 12-nearest-neighbor graph. (B--D) The equal-energy mixture and its true low-frequency and localized packet components. (E--F) The Tight GBDN minus and plus channels for the selected complex-response fit. (G) The target, real learned, and magnitude responses on the graph spectrum. (H) Learned root trajectories and mapped target poles. The inset metrics use the displayed run; per-initialization artifacts were not retained.}
\label{fig:sphere-decomposition}
\end{figure*}

The second mechanism study sweeps four root radii, four root angles, and Chebyshev degrees from 4 to 64. At degree 16, the retained records give rank correlations of $0.873$ between log error and root radius and $-1.000$ between log error and the mapped-target-pole ellipse parameter. The sign reflects faster convergence for a more distant pole. \Cref{fig:pole-distance} reports the full degree sweep and the mapped-pole ordering. All 80 grid records are retained, but the artifact records a composite source hash rather than a source commit and that hash does not match the current implementation. This study is therefore a diagnostic illustration of why radius alone is incomplete, pending a commit-bound rerun.

\begin{figure*}[t]
\centering
\includegraphics[width=0.88\textwidth]{fig_pole_distance.pdf}
\caption{Diagnostic Chebyshev realization study. (A) Maximum symbol error over $[0,2]$ for four radii and four angles. Line style distinguishes angles at a fixed color-coded radius. (B) At degree 16, the Bernstein-ellipse parameter induced by the mapped target pole orders all retained approximation errors monotonically.}
\label{fig:pole-distance}
\end{figure*}

The matched response-efficiency study remains open. It will compare Tight GBDN, Product-sum GBDN, GBDN+, CayleyNet, ChebNetII, and a verified WaveGC implementation at matched parameters and sparse matrix-vector products. Gate C remains partial until those external implementations are verified and run.

\subsection{Preliminary legacy H100 reproduction}
\label{sec:legacy-h100}

The preserved legacy notebook executes on one NVIDIA H100 80GB HBM3, but it
predates Tight GBDN and evaluates only the relaxed GBDN+ model.  All 60
heterophily jobs completed, yet the rerun fails its reproduction criterion:
two primary metrics drift by more than 0.02 and the required run manifest is
absent.  The two Peptides-func jobs retain runtime and peak-memory measurements
but no test predictions for independent metric verification.  We therefore
report the complete protocols and measurements only as a portability
diagnostic in \cref{app:legacy-h100}; they neither validate the local
comparators nor support a performance ordering.

\subsection{Downstream evaluation gate}

Confirmatory downstream claims remain deferred. Heterophily experiments will use all official splits and metrics, at least three training seeds per split, equal hyperparameter-search trials, and confidence intervals that respect the split--seed dependence. LRGB results require a verified graph-level pipeline and official evaluators. The single-run values above do not satisfy this gate and are excluded from the paper's contribution claims.
